\documentclass{article}
\linespread{1.3}
\usepackage[margin=50pt]{geometry}
\usepackage{amsmath, amsthm, amssymb, amsthm, tikz, fancyhdr, graphicx, systeme}
\pagestyle{fancy}
\renewcommand{\headrulewidth}{0pt}
\newcommand{\changefont}{\fontsize{15}{15}\selectfont}

\fancypagestyle{firstpageheader}{
  \fancyhead[R]{
    \changefont
    \parbox[t]{4cm}{ % Adjust width as needed
      Michael Huang\\
      EN.625.633.81\\
      Homework 5
    }
  }
}

\begin{document}

\thispagestyle{firstpageheader}
{\Large 

\section*{1.}

We are given a Laplace distribution with pdf 
\[
  f_X(x) = \frac{\lambda}{2}\exp\{-\lambda|x-\theta|\}
\]
with \(\lambda = 2\) and \(\theta = 2\), for \(-\infty < x < \infty\). To apply the inverse transform method, we first find the CDF. Given the absolute value in the function, we need to find the CDF for two different cases: \\ 
- Case 1: \(x > \theta\)
\[
  F_X(x) = P(X \leq x) = \int_{-\infty}^{x} f_X(x) dx = \int_{-\infty}^{\theta} f_X(x) dx + \int_{\theta}^{x} f_X(x) dx
\]
\[
  = \int_{-\infty}^{\theta} \frac{\lambda}{2}\exp\{-\lambda|x-\theta|\} dx + \int_{\theta}^{x} \frac{\lambda}{2}\exp\{-\lambda|x-\theta|\} dx
\]
\[
  = \int_{-\infty}^{\theta} \frac{\lambda}{2}\exp\{-\lambda(-x-(-\theta))\} dx + \int_{\theta}^{x} \frac{\lambda}{2}\exp\{-\lambda(x-\theta)\} dx
\]
\[
  = \int_{-\infty}^{\theta} \frac{\lambda}{2}\exp\{\lambda(x-\theta)\} dx + \int_{\theta}^{x} \frac{\lambda}{2}\exp\{-\lambda(x-\theta)\} dx
\]
For the first integral, we do substitution with \(u = \lambda(x - \theta)\) (and thus \(\frac{du}{dx} = \lambda\) and \(dx = \frac{1}{\lambda}du\)):
\[
  = \int_{-\infty}^{0} \frac{\lambda}{2}\exp\{u\} \cdot \frac{1}{\lambda} du + \int_{\theta}^{x} \frac{\lambda}{2}\exp\{-\lambda(x-\theta)\} dx 
\]
\[
  = \int_{-\infty}^{0} \frac{1}{2}\exp\{u\} du + \int_{\theta}^{x} \frac{\lambda}{2}\exp\{-\lambda(x-\theta)\} dx 
\]
\[
  = \frac{1}{2}\exp\{u\} |_{-\infty}^{0}  + \int_{\theta}^{x} \frac{\lambda}{2}\exp\{-\lambda(x-\theta)\} dx 
\]
\[
  = \frac{1}{2}\exp\{0\} - 0 + \int_{\theta}^{x} \frac{\lambda}{2}\exp\{-\lambda(x-\theta)\} dx 
\]
\[
  = \frac{1}{2} + \int_{\theta}^{x} \frac{\lambda}{2}\exp\{-\lambda(x-\theta)\} dx 
\]
Do substitution in a similar way for the second integral, with \(u = -\lambda(x - \theta)\), \(\frac{du}{dx} = -\lambda\), and \(dx = -\frac{1}{\lambda}du\):
\[
  = \frac{1}{2} + \int_{0}^{-\lambda(x - \theta)} \frac{\lambda}{2}\exp\{u\} \cdot -\frac{1}{\lambda} du
\]
\[
  = \frac{1}{2} + \int_{0}^{-\lambda(x - \theta)} -\frac{1}{2}\exp\{u\} du
\]
\[
  = \frac{1}{2} - \frac{1}{2}\exp\{u\} |_{0}^{-\lambda(x - \theta)}
\]
\[
  = \frac{1}{2} - \frac{1}{2}(\exp\{-\lambda(x - \theta)\} - \exp\{0\})
\]
\[
  = \frac{1}{2} - \frac{1}{2}(\exp\{-\lambda(x - \theta)\} - 1)
\]
\[
  = \frac{1}{2} + \frac{1}{2} - \frac{1}{2}(\exp\{-\lambda(x - \theta)\})
\]
\[
  = 1 - \frac{1}{2}\exp\{-\lambda(x - \theta)\}
\]
- Case 2: \(x < \theta\)
\[
  F_X(x) = P(X \leq x) = \int_{-\infty}^{x} f_X(x) dx 
\]
\[
  = \int_{-\infty}^{x} \frac{\lambda}{2}\exp\{-\lambda(-x-(-\theta))\} dx
\]
\[
  = \int_{-\infty}^{x} \frac{\lambda}{2}\exp\{\lambda(x-\theta)\} dx
\]
We do substitution as we previously did above:
\[
  = \int_{-\infty}^{\lambda(x-\theta)} \frac{\lambda}{2}\exp\{u\} \cdot \frac{1}{\lambda} du
\]
\[
  = \int_{-\infty}^{\lambda(x-\theta)} \frac{1}{2}\exp\{u\} du
\]
\[
  = \frac{1}{2}\exp\{u\} |_{-\infty}^{\lambda(x-\theta)}
\]
\[
  = \frac{1}{2}(\exp\{\lambda(x-\theta)\} - \exp\{-\infty\})
\]
\[
  = \frac{1}{2}\exp\{\lambda(x-\theta)\}
\]
Now that we have the CDFs, let's invert accordingly. We know that given the CDF has a breakpoint at the value where the cases are equal, which we can evaluate at \(x = \theta\) to be \(\frac{1}{2}\exp\{\lambda(\theta-\theta)\} = \frac{1}{2}\exp\{0\} = \frac{1}{2}\). So we start inverting based on the corresponding values: \\
- Case 1: \(U_i > \frac{1}{2} \) \\
\[
  U_i = 1 - \frac{1}{2}\exp\{-\lambda(x - \theta)\}
\]
\[
  \frac{1}{2}\exp\{-\lambda(x - \theta)\} = 1 - U_i
\]
\[
  \exp\{-\lambda(x - \theta)\} = 2(1 - U_i)
\]
\[
  -\lambda(x - \theta) = \ln\{2(1 - U_i)\}
\]
\[
  x - \theta = -\frac{\ln\{2(1 - U_i)\}}{\lambda}
\]
\[
  x = \theta -\frac{\ln\{2(1 - U_i)\}}{\lambda}
\]
- Case 2: \(U_i < \frac{1}{2} \) \\
\[
  U_i = \frac{1}{2}\exp\{\lambda(x-\theta)\}
\]
\[
  2U_i = \exp\{\lambda(x-\theta)\}
\]
\[
  \ln\{2U_i\} = \lambda(x-\theta)
\]
\[
  \frac{\ln\{2U_i\}}{\lambda} = x - \theta
\]
\[
  x = \frac{\ln\{2U_i\}}{\lambda} + \theta
\]
So to generate the random variables, we do the following: \\
- Generate \(U_1, U_2, \dots, U_n \sim \text{Unif(0,1)}\) \\ 
- For \(U_i > \frac{1}{2}\), set \(X_i = F^{-1}_X(U_i) = \theta - \frac{\ln\{2(1 - U_i)\}}{\lambda} = 2 - \frac{\ln\{2(1 - U_i)\}}{2}\) \\
- For \(U_i < \frac{1}{2}\), set \(X_i = F^{-1}_X(U_i) = \frac{\ln\{2U_i\}}{\lambda} + \theta = 2 + \frac{\ln\{2U_i\}}{2}\) \\
- For \(U_i = \frac{1}{2}\), apply either function above.

}

{\Large 

\section*{2.}

We are given a distribution with pdf 
\[ 
  f_X(x) = \begin{cases} \frac{1}{4} & 0 < x < 1 \\[4pt] x - \frac{3}{4} & 1 \leq x \leq 2 \end{cases} 
\]

\subsection*{(a)} 

As we normally do with the inverse-transform method, let's take the CDF for each section: \\
- \(0 < x < 1\): \\ 
\[
  F_X(x) = \int_0^x f_X(x) dx
\]
\[
  = \int_0^x \frac{1}{4} dx
\]
\[
  = \frac{1}{4}x |_0^x
\]
\[
  = \frac{1}{4}x - 0
\]
\[
  F_X(x) = \frac{x}{4}
\]
- \(1 \leq x \leq 2\): \\
\[
  F_X(x) = \int_0^1 f_X(x) dx + \int_1^x f_X(x) dx
\]
\[
  = \int_0^1 \frac{1}{4} dx + \int_1^x x - \frac{3}{4} dx
\]
\[
  = \frac{1}{4}x |_0^1 + \frac{x^2}{2} - \frac{3}{4}x |_1^x
\]
\[
  = \frac{1}{4}x |_0^1 + [\frac{x^2}{2} - \frac{3}{4}x] |_1^x
\]
\[
  = \frac{1}{4} - 0 + \frac{x^2}{2} - \frac{3}{4}x - \frac{1}{2} + \frac{3}{4}
\]
\[
  = \frac{x^2}{2} - \frac{3}{4}x + \frac{1}{2}
\]
Now, invert accordingly: \\
- \(0 < x < 1\): \\ 
\[
  U_i = \frac{x}{4}
\]
\[
  x = 4U_i
\]
- \(1 \leq x \leq 2\): \\ 
\[
  U_i = \frac{x^2}{2} - \frac{3}{4}x + \frac{1}{2}
\]
\[
  4U_i = 2x^2 - 3x + 2
\]
\[
  0 = 2x^2 - 3x + (2 - 4U_i)
\]
\[
  x = \frac{3 \pm \sqrt{(-3)^2 - 4(2 \cdot (2 - 4U_i))}}{2(2)}
\]
\[
  x = \frac{3 \pm \sqrt{9 - 16 + 32U_i}}{4}
\]
\[
  x = \frac{3 \pm \sqrt{32U_i - 7}}{4}
\]
To see whether we should be using \(+\) or \(-\), let's evaluate at the breakpoint of \(x=1\) and \(U_i = \frac{1}{4}\):
\[
  1 = \frac{3 \pm \sqrt{32(\frac{1}{4})- 7}}{4}
\]
\[
  1 = \frac{3 \pm \sqrt{8- 7}}{4}
\]
\[
  4 = 3 \pm 1
\]
So clearly we should take \(+\) in this case. We would therefore take the following steps to generate the random variables: \\
- Generate \(U_1, U_2, \dots, U_n \sim \text{Unif(0,1)}\) \\ 
- For \(U_i < \frac{1}{4}\), set \(X_i = F^{-1}_X(U_i) = 4U_i\) \\
- For all other values, set \(X_i = F^{-1}_X(U_i) = \frac{3 + \sqrt{32U_i - 7}}{4}\)

\subsection*{(b)}

Given the proposal density \(g(x) = \frac{1}{2} \quad 0 \leq x \leq 2\), for the accept-reject method, we need to find the appropriate constant \( M \geq sup_x[\frac{f(x)}{g(x)}]\). We know that 
\[
  \frac{f(x)}{g(x)} = \frac{f(x)}{\frac{1}{2}} = 2f(x)
\]
We now find the maximum of \(f(x)\). Looking at the function, this is 
\[
  f(2) = 2 - \frac{3}{4} = \frac{5}{4}
\]
Plugging this in, we find that 
\[
  M = 2 \cdot \frac{5}{4} = \frac{5}{2}
\]
The steps to generate the random variable using the accept-reject method are as follows: \\
- Generate \(X_{cand} \sim g(x)\), i.e. \(X \sim \text{Uniform}(0, 2)\) \\
- Generate \(U \sim \text{Unif}(0, 1)\) independently \\
- If \(U \leq \frac{f(X_{cand})}{M \cdot g(x)} = \frac{f(X_{cand})}{\frac{5}{2} \cdot \frac{1}{2}} = \frac{4f(X_{cand})}{5}\), accept \(X_{cand}\)

}

{\Large 

\section*{3.}

\subsection*{(a)} 



\subsection*{(b)}



}

{\Large 

\section*{4.}

We are given the following discrete uniform distribution with pdf:
\[
  f_X(x) = \frac{1}{n+1}
\]
for \(x = 0, 1, 2, \dots, n\), and 0 otherwise. To apply the inverse-transform method to generate a random variable, we first take the CDF:
\[
  F_x(x) = P(X \leq k) = \sum_{k=0}^{x} f_X(x)
\]
\[
  F_x(x) = \sum_{k=0}^{x} \frac{1}{n+1}
\]
\[
  F_x(x) = \frac{x + 1}{n+1}
\]
We would therefore take the following steps to generate samples: \\
- Generate \(U_1, U_2, \dots, U_n \sim \text{Unif(0,1)}\) \\ 
- Set \(X_i = F^{-1}_X (U_i) = (n+1) \cdot  U_i = \lfloor U_i \cdot (n + 1) \rfloor \), with the floor needed due to the discrete nature of the original function. 

}

{\Large 

\section*{5.}



}

\end{document}